\section{Discussion}

The main message of this study is that, within the studied expanded panel, DTI evaluation is benchmark-sensitive in a way that current reporting practice often obscures. A synthetic split can still be useful, but it should not be treated as a universal proxy for real OOD generalization. Once the panel includes pooled-LM baselines, a graph-family probe, and a learner-control classical baseline on the same handcrafted inputs as \texttt{base}, the sign and magnitude of an apparent gain can still change when the benchmark changes.

An important design choice in this paper is the use of a deliberately expanded but still finite external model panel. A larger model zoo could be added later, but the present panel is already broad enough to invalidate the narrow objection that the earlier ranking changes were artifacts of comparing only pooled-LM variants. In particular, \texttt{RF} and \texttt{base} share the same handcrafted inputs but use different learners, while \texttt{GraphDTA-style} adds a graph-family inductive bias. The new results therefore support a stronger statement than in the smaller panel: benchmark-qualified ranking changes persist even when the comparison spans classical, fixed-feature neural, pooled-LM, and graph-based systems. At the same time, the present paper should still not be read as an exhaustive field-wide census.

The paper also has clear limitations. Although the support benchmarks \texttt{BindingDB\_Ki} and \texttt{DAVIS} remain useful counterexamples, they are still supporting rather than co-equal headline axes. The expanded model family is much stronger than the original in-repo probe panel, but it is still partial rather than exhaustive. We now add two targeted robustness probes, and both make the manuscript more disciplined. The non-patent temporal benchmark shows that the temporal ordering is not confined to patent provenance alone, while the scaffold-drug probe shows that the original synthetic \texttt{unseen\_drug} split was chemically permissive enough that the earlier \texttt{HyperPCM} edge does not survive a stricter holdout. Even with these additions, the patent benchmark is still temporal/provenance OOD rather than a pure time-only shift, so the current paper should not claim that time alone explains the observed ranking change. A further caveat is that ChemBERTa and ESM2 were pretrained on broad corpora that may already include molecular and protein entities appearing in our temporal test sets. We therefore cannot rule out that part of the earlier \texttt{DTI-LM} edge on temporal benchmarks reflects broader pretraining coverage of chemical or protein space rather than temporal robustness per se. At the same time, the expanded panel now shows that the main temporal winner is no longer a pooled-LM model but a classical learner-control baseline using the same handcrafted representation as \texttt{base}. This weakens any mechanistic story that rests only on LM pretraining and shifts the more credible conclusion toward benchmark sensitivity and learner sensitivity. We also do not evaluate downstream biological decision quality. This paper should therefore be read as a well-supported benchmark case study rather than an exhaustive census of all real-OOD DTI settings. These limitations should be stated plainly, but they do not weaken the more disciplined central finding that synthetic and real-OOD rankings can disagree sharply in the tested expanded panel.
